\section{Literature Review}\label{sec:lit}

\subsection{Foundations of volatility spillovers}

The idea that shocks propagate across financial assets has long been central to the study
of interdependence and contagion. Early empirical work mainly relied on cross-market
correlations and cointegration to infer common dynamics. However, as
\citet{ForbesRigobon2002} showed, correlations tend to rise mechanically during turbulent
periods even without genuine contagion, leading to potential misinterpretation. This
limitation prompted the development of methodologies capable of capturing both the
magnitude and direction of volatility transmission. \citet{Yilmaz2006} provided one of
the first applications of variance-decomposition techniques to measure volatility
spillovers among European equity markets, laying the foundation for the connectedness
framework later formalized by \citet{DieboldYilmaz2009}. This shift from static
correlation analysis to variance-decomposition models marked a new, directional
understanding of financial interdependence.

\subsection{The Diebold--Yilmaz connectedness framework}

\citet{DieboldYilmaz2009} introduced a unified method to measure spillovers through the
share of each variable's forecast error variance explained by innovations in others within
a VAR system. This decomposition separates total forecast uncertainty into own and cross
components, yielding a systematic measure of interdependence. \citet{DieboldYilmaz2012}
extended the approach by adopting the Generalized Forecast Error Variance Decomposition
(GFEVD) of \citet{KoopPesaranShin1996} and \citet{PesaranShin1998}, which eliminates
sensitivity to variable ordering. Subsequent work incorporated directional
(\emph{to}/\emph{from}) and net indices \citep{DieboldYilmaz2014}, as well as dynamic
versions based on rolling estimation \citep{DieboldYilmaz2015}. These developments
transformed volatility spillover analysis into a flexible, dynamic, and interpretable
framework that can be consistently applied across markets and time periods.

\subsection{Connectedness as a network of volatility transmission}

The forecast error variance decomposition matrix can be viewed as a weighted, directed
network where nodes represent financial institutions and edges quantify volatility
transmission. This perspective enables analysis through network theory, linking
econometric measures to the topology of systemic risk. Metrics such as in-degree,
out-degree, and overall density capture how shocks are propagated and absorbed across
institutions. The Total Connectedness Index (TCI) summarizes system-wide
interconnectedness, while directional and net measures identify dominant transmitters and
receivers of volatility. Examining how clustering, centrality, and network density evolve
across regimes provides deeper insight into contagion channels and the concentration of
systemic influence \citep{DieboldYilmaz2015}.

\subsection{Volatility measures and data frequency}

Accurate volatility estimation is fundamental to reliable connectedness measurement.
Range-based estimators \citep{AlizadehBrandtDiebold2002} and realized volatility from
high-frequency returns \citep{AndersenBollerslevDieboldLabys2003} provide more efficient
and less noisy proxies than squared returns. Because they respond quickly to information
arrivals, these measures are ideal for detecting rapid volatility transmission.
Equity-based volatilities are particularly informative for financial institutions, as
stock prices incorporate forward-looking market perceptions of solvency and liquidity
conditions, making them suitable for monitoring systemic stress in real time.

\subsection{Empirical evidence for Europe}

Applications of the Diebold--Yilmaz framework to European and trans-Atlantic data reveal
strong time variation in spillovers. \citet{DieboldYilmaz2015TA} show that during the
GFC, volatility mainly flowed from the U.S.\ to Europe, became bidirectional after the
Lehman collapse, and reversed during the sovereign debt crisis as stress in European banks
spilled back to U.S.\ markets. Within Europe, heterogeneity across banks, insurers, and
diversified financials reflects differences in funding models, sovereign exposures, and
business lines. \citet{BlackCorreaHuangZhou2016} demonstrate that systemic risk among
European banks increased sharply during the sovereign crisis, while
\citet{BonaldiHortacsuKastl2015} identify repo market frictions and collateral channels
as key propagation mechanisms. Together, these studies show that connectedness is
regime-dependent, rising sharply during crises and declining in calmer periods.

\subsection{Systemic-risk measures and early-warning indicators}

The connectedness approach complements other systemic-risk methodologies developed in
parallel. CoVaR \citep{AdrianBrunnermeier2016} captures tail co-movement between
institutions, MES and SES \citep{Acharya2017} quantify expected capital shortfalls under
systemwide distress, and SRISK \citep{BrownleesEngle2017} estimates the capital each
institution would require to remain solvent in adverse scenarios. Complementary indicators
such as the Absorption Ratio \citep{KritzmanLiPageRigobon2011} and Granger-causality
networks \citep{BillioEtAl2012JFE} serve as early-warning signals by revealing increased
commonality and directional dependence before crises. Although these measures differ
conceptually, they all converge on a key insight: systemic fragility arises when risk
transmission becomes concentrated and market co-movement intensifies.

\subsection{Contagion versus interdependence}

A persistent debate distinguishes true contagion (excess transmission beyond what
fundamentals explain) from normal interdependence. The connectedness framework does not
claim structural identification but provides a dynamic map of predictive linkages.
Monitoring time variation in connectedness helps detect phases where co-movement exceeds
historical norms, consistent with contagion episodes. When combined with macro-financial
markers such as liquidity stress or policy events, these measures can differentiate
between endogenous and exogenous drivers of systemic shocks.

\subsection{Economic mechanisms of transmission}

Economic theory helps interpret why connectedness fluctuates.
\citet{BrunnermeierPedersen2009} model funding and market liquidity spirals that
propagate price shocks across assets, while \citet{DanielssonShinZigrand2012} show how
procyclical leverage amplifies volatility through balance-sheet feedbacks.
\citet{KingMaier2009} emphasize how hedge funds and prime brokers transmit stress via
margin and funding constraints. These mechanisms explain why connectedness tends to surge
during liquidity shortages and deleveraging episodes, underscoring its link to real
financial frictions.

\subsection{Correlation-based networks and minimum spanning trees}

In parallel with the econometric spillover literature, a large econophysics and
data-driven strand models financial markets directly through empirical dependence
structures. Starting from \citet{Mantegna1999}, this work constructs correlation-based
distance matrices, filters them into minimum spanning trees (MSTs) and related sparse
graphs, and analyzes the resulting topology (clusters, hierarchy, and
centrality) to understand market structure and diversification. The survey of
\citet{MartiEtAl2020} reviews two decades of research on correlations, hierarchies,
networks, and clustering in financial markets, including planar filtered graphs, community
detection, and their applications to portfolio construction. These methods are largely
statistical and model-free: they do not impose a specific dynamic parametric structure,
but instead infer network structure from empirical correlations (or more general
dependence measures) and use it to design robust, diversification-oriented portfolio
rules. This provides a natural benchmark against which to compare model-based
connectedness indices.

\subsection{Multivariate GARCH: DCC versus BEKK}

Although VAR--GFEVD connectedness remains the benchmark, covariance-based multivariate
GARCH models offer a complementary view of time-varying dependence. The BEKK
specification of \citet{EngleKroner1995} guarantees positive definiteness of the
conditional covariance matrix and allows for rich cross-dynamics between variances and
covariances. However, even the ``diagonal'' BEKK model implies $O(N^2)$ parameters in
an $N$-dimensional system, and the full BEKK quickly becomes over-parameterized as $N$
grows. This high dimensionality makes estimation computationally demanding and can lead
to convergence issues or unstable coefficients, especially when models are re-estimated
on rolling windows.

By contrast, the Dynamic Conditional Correlation (DCC) model of \citet{Engle2002DCC}
separates univariate volatility dynamics from the evolution of the correlation matrix.
Each series is first modeled by a univariate GARCH, and a small number of additional
parameters governs the dynamics of the common correlation process. As a result, the
number of DCC parameters does not increase with $N$, making it far more tractable in
high-dimensional systems of financial institutions. DCC also delivers a directly
interpretable sequence of correlation matrices that can be mapped into time-varying
network structures, consistent with the connectedness perspective. For empirical
applications like ours, where $N$ is relatively large and dynamic analysis relies on
rolling estimation, DCC--GARCH thus offers a practical and robust benchmark, while BEKK
remains conceptually useful but less feasible in routine implementation.

\subsection{Multilayer and clustered network connectedness}

Recent work emphasizes that systemic risk should be analyzed jointly through variance,
covariance, and network perspectives, often in a multilayer setting. Multilayer network
models distinguish linkages by sector, maturity, or volatility proxy, revealing how
cross-layer interactions reinforce systemic risk \citep{WangEtAl2021}. Building on the
Diebold--Yilmaz framework, \citet{BuchwalterDieboldYilmaz2025} propose a clustered
network connectedness approach, in which rolling connectedness matrices are summarized
into a small number of latent regimes. By clustering the time-varying spillover structure,
they identify periods with distinct patterns of cross-market transmission and show how
these regimes map to global equity-market events. This line of research motivates our use
of clustering techniques on rolling connectedness measures to uncover volatility regimes
in European financial institutions and to relate them to macro-financial stress episodes.

\subsection{Frequency-domain and dynamic extensions}

\citet{BarunikKrehlik2018} extend the connectedness framework to the frequency domain,
decomposing spillovers into short-term and long-term components. Short-frequency
connectedness often reflects transient shocks linked to liquidity or sentiment, whereas
long-frequency connectedness captures persistent macro-financial co-movement. Combining
frequency decomposition with rolling estimation provides a comprehensive time--frequency
view of systemic interdependence.

\subsection{Synthesis and research gap}

The literature consistently shows that volatility spillovers are time-varying, intensify
during crises, and reveal key propagation channels of systemic risk. Connectedness
measures provide a transparent, network-consistent way to map these linkages, while DCC,
BEKK, multilayer, clustered, and correlation-network frameworks offer alternative lenses
on dependence dynamics. Yet few studies have jointly applied the full Diebold--Yilmaz framework to European
financial institutions while (i)~formally testing for structural breaks in the
connectedness process, (ii)~decomposing spillovers across frequency bands to separate
transient contagion from durable structural integration, (iii)~identifying connectedness
regimes directly from GFEVD feature vectors without pre-labelling crisis episodes, and
(iv)~benchmarking model-based spillovers against model-free correlation and MST network
centrality across the full institutional cross-section. Addressing this gap, the present
paper delivers each of these analyses on a broad panel of 63 European financial
institutions, providing both formal statistical tests and cross-method validation of the
identified systemic risk rankings.
